\documentclass[UTF8,11pt,a4paper]{ctexart}
\usepackage{amsmath,amssymb,geometry,enumitem,hyperref}
\geometry{margin=2cm}\setlist{itemsep=2pt,topsep=4pt}
\hypersetup{colorlinks=true,linkcolor=blue}
\newcommand{\BOTH}[2]{\par\noindent\textbf{中文。}\ #1\par\noindent\textbf{English.}\ #2\par}
\newcommand{\EX}[2]{\par\noindent\textbf{例子。}\ #1\par\noindent\textbf{Example.}\ #2\par}
\newcommand{\READ}[1]{\par\noindent\textbf{公式英文读法 / How to read it.} #1\par}
\newcommand{\ASK}[4]{\begin{samepage}\par\noindent\textbf{可能问题。}\ #1\par\noindent\textbf{Possible question.}\ #2\par\noindent\textbf{回答要点。}\ #3\par\noindent\textbf{Key answer.}\ #4\medskip\end{samepage}}
\title{数学分析（三）APS 面试准备\\[4pt]\large 实数完备性、级数与多变量极限}
\author{Mathematical Analysis III: Bilingual Oral Review}
\date{}
\begin{document}\maketitle

{\footnotesize\noindent\textbf{范围 / Scope.} 依据“数学分析”文件夹中的讲义及第三学期范围：第 2 章 2.4、第 6 章 6.1--6.3、第 11 章 11.1--11.4；带星号内容及未列章节均排除。 Based on the local lecture notes and the third-semester outline: Section 2.4, Sections 6.1--6.3, and Sections 11.1--11.4. Starred and unlisted topics are excluded.\par}

\section*{课程概述 | Course overview}
\BOTH{我在第三学期围绕一个基础问题展开：为什么实数上的极限过程不会“缺点”？确界原理、单调收敛、闭区间套、Bolzano--Weierstrass、Cauchy 收敛和 Heine--Borel 定理，是实数完备性的不同表达。级数部分研究无穷求和何时有确定值、不同判别法适用于什么条件；多变量极限则说明在平面或高维空间中，趋近点有无数方向，必须控制所有方向而不能只检查几条曲线。}{In my third semester, I studied a basic question: why do limits over the real numbers not develop gaps? The supremum property, monotone convergence, nested intervals, Bolzano--Weierstrass, the Cauchy criterion, and Heine--Borel are different formulations of completeness. I also studied when infinite sums have definite values, which convergence tests apply, and why multivariable limits require control along every direction of approach.}

\BOTH{这些思想有实际计算意义。判定一个数值迭代是否收敛，先证明它是 Cauchy 或单调有界；估计无穷过程的总贡献，要选适合项型的级数判别；验证二维函数极限，沿 $y=0$ 检查不够，必须控制任意 $(x,y)$ 同时趋近。}{These ideas matter in computation. To justify an iterative limit, one may prove the sequence is Cauchy or monotone and bounded. To estimate an infinite accumulated contribution, choose a series test matching the terms. For a two-variable limit, checking only $y=0$ is insufficient; all approaches must be controlled.}

\section{实数完备性 | Completeness of the real numbers}
\subsection{上确界与六种等价表述}
\BOTH{集合 $A\subset\mathbb R$ 的上界 $M$ 满足 $x\le M$ 对所有 $x\in A$；若上界中最小者存在，称为上确界 $\sup A$。上确界不一定属于 $A$，例如 $(0,1)$ 的上确界是 $1$。实数完备性说每个非空且有上界的实数集合都有上确界。}{An upper bound $M$ of $A\subset\mathbb R$ satisfies $x\le M$ for every $x\in A$. The least upper bound, if it exists, is $\sup A$ and need not belong to $A$; for example, $\sup(0,1)=1$. Completeness says every nonempty real set bounded above has a supremum.}
\BOTH{在实数系中，以下性质等价：确界原理；单调有界数列收敛；长度趋零的嵌套非空闭区间有唯一公共点；每个有界数列有收敛子列；每个 Cauchy 数列收敛；每个闭有界集合紧致（每个开覆盖有有限子覆盖）。它们是同一完备结构的不同说法，不表示在任意集合或任意度量空间中都自动等价。}{For the real numbers, these properties are equivalent: the least-upper-bound property; convergence of every bounded monotone sequence; the unique common point of nested nonempty closed intervals whose lengths tend to zero; a convergent subsequence in every bounded sequence; convergence of every Cauchy sequence; and compactness of every closed bounded set, meaning every open cover has a finite subcover. These are formulations of the same completeness structure, not automatic equivalences in arbitrary spaces.}
\BOTH{单调收敛定理直接由确界得到：递增有界数列的项集有上确界 $L$；按上确界定义，任意 $\varepsilon>0$ 都有某项大于 $L-\varepsilon$，之后所有项夹在 $L-\varepsilon$ 与 $L$ 之间，所以趋于 $L$。反方向的等价性需要构造相应的单调列或区间列。}{The monotone convergence theorem follows directly from the supremum. Let $L$ be the supremum of an increasing bounded sequence. For every $\varepsilon>0$, some term exceeds $L-\varepsilon$; all later terms then lie between $L-\varepsilon$ and $L$, so the sequence converges to $L$. The reverse equivalences use suitable constructions of monotone sequences or intervals.}
\EX{令 $a_n=1-1/n$。它递增且以上界 $1$ 控制；其项集上确界为 $1$，因此 $a_n\to1$，虽然每一项都小于 $1$。这说明极限可以是数列从未取到的上确界。}{The sequence $a_n=1-1/n$ increases and is bounded above by one. Its supremum is one, so it converges to one even though no term equals one. A limit can be a supremum never attained by the sequence.}
\READ{“The supremum of A is the least upper bound of A.”}

\subsection{闭区间套、Bolzano--Weierstrass 与 Cauchy}
\BOTH{闭区间套定理：若 $I_1\supseteq I_2\supseteq\cdots$ 是非空闭区间且长度趋于零，则交集恰有一个点。证明思路是左端点递增、右端点递减，利用实数完备性得到共同极限；区间长度趋零保证唯一。}{The nested interval theorem states that if $I_1\supseteq I_2\supseteq\cdots$ are nonempty closed intervals whose lengths tend to zero, their intersection contains exactly one point. The left endpoints increase and the right endpoints decrease; completeness gives a common limiting point, and shrinking lengths give uniqueness.}
\EX{二分区间 $[0,1]$：每一步保留包含给定实数 $x$ 的一半，得到长度 $2^{-n}$ 的闭区间套。唯一交点就是 $x$。这也是二进制逼近的数学基础。}{Repeatedly bisect $[0,1]$ and keep the half containing a fixed real number $x$. The resulting intervals have lengths $2^{-n}$ and their unique common point is $x$, providing the basis for binary approximation.}
\BOTH{Bolzano--Weierstrass 定理说有界实数列必有收敛子列。证明可不断在区间中选取含有无穷多项的半区间，构造闭区间套，再从中选指标递增的项。Heine--Borel 定理在 $\mathbb R^n$ 中说集合紧致当且仅当闭且有界；在 $\mathbb R$ 中它把有限覆盖性质与有界闭集联系起来。}{Bolzano--Weierstrass says every bounded real sequence has a convergent subsequence. One proof repeatedly selects a half-interval containing infinitely many terms, constructs nested intervals, and chooses increasing indices. In $\mathbb R^n$, Heine--Borel says a set is compact exactly when it is closed and bounded.}
\BOTH{Cauchy 准则只比较数列内部的项：对任意 $\varepsilon>0$，当 $m,n$ 都足够大时 $|a_m-a_n|<\varepsilon$。它不需要预先知道候选极限；在 $\mathbb R$ 中完备性保证每个 Cauchy 列确有实数极限。数值计算里可先用 Cauchy 估计验证逼近稳定，再识别极限的具体值。}{The Cauchy criterion compares terms within the sequence: for every $\varepsilon>0$, all sufficiently late pairs satisfy $|a_m-a_n|<\varepsilon$. It does not require knowing the candidate limit. Completeness of $\mathbb R$ ensures that every Cauchy sequence has a real limit. In numerical work, one can establish stability of an approximation before identifying its exact limit.}
\EX{若 $a_n=\sum_{k=1}^n2^{-k}$，则 $m>n$ 时 $|a_m-a_n|=\sum_{k=n+1}^m2^{-k}<2^{-n}$。因此部分和是 Cauchy 数列，并收敛到 $1$。}{For $a_n=\sum_{k=1}^n2^{-k}$ and $m>n$, the difference is less than $2^{-n}$. Thus the partial sums are Cauchy and converge to one.}

\EX{开区间 $(0,1)$ 的开覆盖 $U_n=(1/n,1)$（$n\ge2$）覆盖整个区间，但任意有限多个中最大的指标 $N$ 仍遗漏 $(0,1/N]$ 中的点，因此没有有限子覆盖。这个例子直接显示非紧致集合可能需要无限多个覆盖集。}{The sets $U_n=(1/n,1)$ cover $(0,1)$, but every finite subcollection has a largest index and misses points sufficiently close to zero. This open cover has no finite subcover.}
\BOTH{Bolzano--Weierstrass 也可用于证明连续函数在紧集上取得最大值：先取函数值趋近上确界的点列，再取收敛子列；闭性保证极限点仍在集合中，连续性保证该点函数值就是上确界。}{Bolzano--Weierstrass also helps prove that a continuous function attains its maximum on a compact set. Choose a sequence whose function values approach the supremum and extract a convergent subsequence. Closedness keeps the limit in the set, and continuity makes its value equal to the supremum.}
\ASK{Cauchy 准则和单调有界定理分别适合什么情形？}{When are the Cauchy criterion and monotone convergence theorem useful?}{若能直接比较相邻尾项或尾和，Cauchy 准则合适；若递推关系容易证明单调和上下界，单调有界定理通常更短。二者都依赖实数完备性来保证极限存在。}{Use the Cauchy criterion when tail differences can be bounded directly. Use monotone convergence when a recurrence makes monotonicity and bounds easy to prove. Both rely on completeness to guarantee a limit.}

\section{第六章：数项级数 | Chapter 6: Numerical series}
\subsection{部分和与 Cauchy 收敛准则}
\BOTH{级数 $\sum a_n$ 收敛，定义为部分和 $s_N=\sum_{n=1}^Na_n$ 收敛。Cauchy 级数准则要求：对每个 $\varepsilon>0$，存在 $N$，使任意 $m>n\ge N$ 都有 $|a_{n+1}+\cdots+a_m|<\varepsilon$。它检验尾和是否整体变小。}{The series $\sum a_n$ converges when its partial sums $s_N=\sum_{n=1}^Na_n$ converge. The Cauchy criterion requires every sufficiently late finite tail to have absolute value below any prescribed $\varepsilon$. It tests whether all tails become small.}
\BOTH{必要条件是 $a_n\to0$，因为 $a_n=s_n-s_{n-1}$；但它不充分。几何级数和 $p$ 级数是基准：$\sum r^n$ 在 $|r|<1$ 收敛；$\sum1/n^p$ 当且仅当 $p>1$ 收敛。}{A necessary condition is $a_n\to0$, since $a_n=s_n-s_{n-1}$, but this is not sufficient. Geometric series converge for $|r|<1$, and $p$-series $\sum1/n^p$ converge exactly when $p>1$.}
\EX{调和级数通项 $1/n\to0$，但按二进分组，每组 $2^{k-1}<n\le2^k$ 的和至少为 $1/2$，因此部分和无界，级数发散。}{The harmonic terms tend to zero, but grouping indices $2^{k-1}<n\le2^k$ gives a sum at least $1/2$ in every block. The partial sums are unbounded, so the series diverges.}

\subsection{正项级数判别法}
\BOTH{正项级数部分和单调递增，因此它收敛当且仅当部分和有上界。比较判别：若 $0\le a_n\le b_n$ 且 $\sum b_n$ 收敛，则 $\sum a_n$ 收敛；若 $a_n\ge b_n\ge0$ 且 $\sum b_n$ 发散，则 $\sum a_n$ 发散。极限比较在 $a_n/b_n\to c\in(0,\infty)$ 时给出相同敛散性。}{For a positive-term series, partial sums increase, so convergence is equivalent to boundedness of the partial sums. The comparison test transfers convergence downward and divergence upward. Limit comparison gives the same behavior when $a_n/b_n$ tends to a positive finite constant.}
\EX{比较 $\sum_{n\ge1}1/(n^2+1)$ 与 $\sum1/n^2$：有 $1/(n^2+1)\le1/n^2$，后者收敛，故原级数收敛。比较时必须保证项非负并方向正确。}{Since $1/(n^2+1)\le1/n^2$ and $\sum1/n^2$ converges, comparison proves convergence. The terms must be nonnegative and the inequality must point in the correct direction.}
\BOTH{积分判别适用于 $a_n=f(n)$，其中 $f$ 在 $[1,\infty)$ 连续、正且递减：级数与 $\int_1^\infty f(x)dx$ 同敛散。对正项级数，若相邻项比的极限小于 $1$ 则比值判别给出收敛，大于 $1$ 则给出发散；根值判别令 $L=\limsup\sqrt[n]{a_n}$，$L<1$ 时收敛，$L>1$ 时发散。两种判别在临界值 $1$ 时都无结论。}{The integral test applies to $a_n=f(n)$ for a positive, continuous, decreasing function: the series and improper integral have the same convergence behavior. For a positive-term series, a ratio limit below one proves convergence and a limit above one proves divergence. For the root test, let $L=\limsup\sqrt[n]{a_n}$; $L<1$ proves convergence and $L>1$ proves divergence. Both tests are inconclusive at the critical value one.}
\EX{$\sum1/n^p$ 用积分判别：$\int_1^\infty x^{-p}dx$ 在 $p>1$ 收敛，故级数也恰在 $p>1$ 收敛。对 $a_n=1/n$，相邻项比趋于 $1$，比值判别无结论，不能误判为收敛。}{The integral test gives convergence of $\sum1/n^p$ exactly for $p>1$. For $1/n$, the ratio of consecutive terms tends to one, so the ratio test is inconclusive; it does not imply convergence.}
\BOTH{当 $a_n>0$ 且 $a_n/a_{n+1}\to1$ 时，仅凭这个极限不能判敛散。例 $1/n$ 对应调和级数发散，$1/n^2$ 对应收敛级数，两者相邻项比都趋于 $1$。}{Knowing only that $a_n/a_{n+1}\to1$ does not determine convergence. Both $1/n$ and $1/n^2$ have this ratio limit, yet one series diverges and the other converges.}

\BOTH{若 $a_n,b_n>0$ 且 $a_n/b_n\to c>0$，则从某项起二者相差至多固定倍数，所以敛散性相同。积分判别的直觉是用递减函数下方或上方的矩形面积夹住积分与级数的部分和。}{If $a_n,b_n>0$ and $a_n/b_n\to c>0$, then eventually the terms differ by at most fixed factors, so the series have the same convergence behavior. The integral test uses rectangles under or above a decreasing graph to compare the integral with partial sums.}

\subsection{任意项级数、绝对收敛与判别法}
\BOTH{若 $\sum|a_n|$ 收敛，称 $\sum a_n$ 绝对收敛；它必收敛，且任意重排不改变和。若原级数收敛但绝对值级数发散，称条件收敛。Leibniz 判别：$a_n\ge0$ 单调递减且趋于零时，$\sum(-1)^{n-1}a_n$ 收敛，余项绝对值不超过首个遗漏项。}{A series is absolutely convergent if $\sum|a_n|$ converges; then the original series converges and every rearrangement has the same sum. Conditional convergence means the original series converges but the absolute-value series diverges. Leibniz's test says an alternating series converges when its positive terms decrease to zero; the error is at most the first omitted term.}
\EX{$1-1/2+1/3-1/4+\cdots$ 按 Leibniz 判别收敛，但绝对值级数是调和级数而发散，所以它条件收敛。取前 $N$ 项时误差不超过 $1/(N+1)$。}{The alternating harmonic series converges by Leibniz's test, but its absolute-value series is harmonic and diverges. It is conditionally convergent, and the error after $N$ terms is at most $1/(N+1)$.}
\BOTH{Dirichlet 判别：若 $A_n=\sum_{k=1}^na_k$ 有界，且 $b_n$ 单调趋于零，则 $\sum a_nb_n$ 收敛。Abel 判别：若 $\sum a_n$ 收敛、$b_n$ 单调且有界，则 $\sum a_nb_n$ 收敛。前者控制振荡因子的部分和，后者利用已收敛级数的分部求和。}{Dirichlet's test states that $\sum a_nb_n$ converges if the partial sums of $a_n$ are bounded and $b_n$ decreases monotonically to zero. Abel's test states convergence when $\sum a_n$ converges and $b_n$ is monotone and bounded. The first controls oscillatory partial sums; the second follows from summation by parts.}
\EX{取 $a_n=(-1)^{n-1}$、$b_n=1/n$，$a_n$ 的部分和只在 $0,1$ 间取值，$b_n$ 递减趋零，Dirichlet 判别给出交错调和级数收敛。}{For $a_n=(-1)^{n-1}$, the partial sums are bounded; $b_n=1/n$ decreases to zero. Dirichlet's test proves convergence of the alternating harmonic series.}
\BOTH{若两个级数绝对收敛，其 Cauchy 乘积系数 $c_n=\sum_{k=0}^na_kb_{n-k}$ 满足 $\sum c_n=(\sum a_n)(\sum b_n)$。绝对收敛也允许重排并保持和；条件收敛级数则可能通过重排改变和，甚至重排成任意指定的实数。这说明“无穷和”必须连同收敛类型讨论。}{For two absolutely convergent series, their Cauchy product has coefficients $c_n=\sum_{k=0}^na_kb_{n-k}$ and sum equal to the product of the original sums. Absolute convergence permits rearrangement without changing the sum. A conditionally convergent series can be rearranged to a different value, even any prescribed real number. Thus an infinite sum must be discussed together with its convergence type.}
\ASK{为什么绝对收敛比条件收敛更稳定？}{Why is absolute convergence more robust than conditional convergence?}{绝对值级数收敛时，尾项绝对值总和可任意小，因此重排不会改变总和；条件收敛依赖正负项的抵消，重排会改变抵消顺序和结果。}{When the absolute-value series converges, the total magnitude of the tail is arbitrarily small, so rearranging terms cannot change the sum. Conditional convergence depends on cancellation between positive and negative terms, and rearrangement can change that cancellation.}

\BOTH{绝对收敛保证重排稳定，因为尾项绝对值总和可以任意小，任何排列都无法从小尾项中制造有限的额外贡献。条件收敛则依赖正负项抵消。交错调和级数按原顺序的和为 $\log2$，重排却可改变和；所以应说明级数是绝对还是条件收敛。}{Absolute convergence makes rearrangement stable because the total magnitude of the tail can be made arbitrarily small. Conditional convergence depends on cancellation between positive and negative terms. The alternating harmonic series sums to $\log2$ in its usual order, but rearrangement can change the value. State whether convergence is absolute or conditional.}
\ASK{比值判别的极限等于 $1$ 能说明什么？}{What does a ratio-test limit of one tell us?}{它不能决定敛散性。调和级数与 $p=2$ 的级数相邻项比都趋于 $1$，前者发散、后者收敛；应改用其他判别法。}{It is inconclusive. Both the harmonic series and the $p$-series with $p=2$ have ratio limit one, but one diverges and the other converges. Use another test.}

\section{第十一章：多变量极限 | Chapter 11: Limits in several variables}
\subsection{\texorpdfstring{$\mathbb R^n$}{R n} 的距离、集合与紧致性}
\BOTH{点 $x=(x_1,\ldots,x_n)$ 的 Euclidean 范数为 $\|x\|=\sqrt{x_1^2+\cdots+x_n^2}$，距离为 $d(x,y)=\|x-y\|$。序列 $x_k\to x$ 等价于 $\|x_k-x\|\to0$，也等价于每个坐标分别收敛。开球 $B(x,r)=\{y:\|y-x\|<r\}$ 是距离语言中的邻域。}{The Euclidean norm is $\|x\|=\sqrt{x_1^2+\cdots+x_n^2}$ and the distance is $d(x,y)=\|x-y\|$. A sequence converges in norm exactly when each coordinate converges. The open ball $B(x,r)$ is the neighborhood defined by distance.}
\BOTH{集合开，若每个点周围都有完全包含在集合内的小球；集合闭，若它包含所有收敛序列的极限；有界，若全体点到某固定点的距离有统一上界。紧致可用开覆盖定义，也可在有限维 Euclidean 空间用“闭且有界”刻画。道路连通指任意两点都能由集合内连续路径连接；它蕴含连通，但不应与紧致混淆。}{A set is open if every point has a small ball contained in it, closed if it contains limits of its convergent sequences, and bounded if all points stay within a fixed distance of some point. Compactness is defined by finite subcovers and, in finite-dimensional Euclidean space, is equivalent to closedness and boundedness. Path connectedness means any two points can be joined by a continuous path inside the set; it implies connectedness but is distinct from compactness.}
\EX{闭圆盘 $\{(x,y):x^2+y^2\le1\}$ 在 $\mathbb R^2$ 中闭且有界，所以紧致；去掉边界的开圆盘仍有界但不闭，因此不紧致。}{The closed unit disk is closed and bounded, hence compact. The open unit disk is bounded but not closed, so it is not compact.}

\subsection{多元极限与累次极限}
\BOTH{若 $a$ 是定义域的聚点，$\lim_{x\to a}f(x)=L$ 表示对每个 $\varepsilon>0$ 存在 $\delta>0$，使 $0<\|x-a\|<\delta$ 时 $|f(x)-L|<\varepsilon$。关键是一个 $\delta$ 同时控制所有方向。多元 Heine 判据等价地要求每个趋向 $a$ 的点列都使函数值趋向 $L$。}{A multivariable limit equals $L$ if every domain point sufficiently close to $a$, in Euclidean distance and distinct from $a$, has function value within $\varepsilon$ of $L$. A single $\delta$ must control every direction. Equivalently, every sequence approaching $a$ must have function values approaching $L$.}
\EX{考察 $f(x,y)=xy/(x^2+y^2)$，原点处不定义。沿 $y=0$ 趋近时值恒为 $0$；沿 $y=x$ 趋近时值恒为 $1/2$。两条路径给出不同值，所以原点极限不存在。}{For $f(x,y)=xy/(x^2+y^2)$, the value along $y=0$ is zero, while along $y=x$ it is $1/2$. Since the path limits differ, the limit at the origin does not exist.}
\BOTH{沿任意一条曲线得到相同结果仍不足以单独证明多元极限；证明通常要给出与方向无关的估计，例如 $|f(x)-L|\le C\|x-a\|$。累次极限先固定一个变量再令另一个趋近。累次极限存在且相等不必推出联合极限存在，上述函数的两个累次极限都是 $0$，但联合极限不存在。}{Checking one path can disprove a multivariable limit, but agreement along a few paths does not prove it. A proof usually gives a direction-independent bound such as $|f(x)-L|\le C\|x-a\|$. Iterated limits take one variable at a time. Even equal iterated limits do not imply a joint limit: both iterated limits of the preceding example are zero, while its joint limit fails to exist.}
\READ{“The limit of f of x as x approaches a equals L in Euclidean norm.”}

\subsection{连续性、紧致性与一致连续}
\BOTH{多元函数 $f$ 在 $a$ 连续，若 $f(a)$ 有定义且极限等于函数值。连续函数的和、积、复合保持连续；向量值函数连续当且仅当各分量连续。紧致集上的连续实值函数有界并取得最大最小值；若定义域道路连通，则连续函数的像具有相应的介值性质。}{A multivariable function is continuous at $a$ if its limit equals its value there. Sums, products, and compositions of continuous functions remain continuous; a vector-valued map is continuous exactly when each component is. A continuous real-valued function on a compact set is bounded and attains its extrema. The continuous image of a path-connected set has the corresponding intermediate-value property.}
\EX{温度场 $T(x,y)=100-x^2-y^2$ 在闭单位圆盘连续，因此一定取得最高、最低温度。直接计算得中心最高为 $100$，边界最低为 $99$。紧致性保证极值存在，计算函数才给出具体位置和值。}{The temperature field $T(x,y)=100-x^2-y^2$ is continuous on the closed unit disk, so it attains its maximum and minimum. Direct calculation gives a maximum of $100$ at the center and a boundary minimum of $99$. Compactness guarantees existence; calculation identifies the values.}
\BOTH{在紧致集合上连续函数一致连续：对所有点共享同一个 $\delta$。这一结论对近似误差控制很重要，因为它避免不同位置需要越来越小的步长。非紧致区域上连续不保证一致连续，例如 $f(x)=x^2$ 在整个实线上不一致连续；但在每个闭有界区间上一致连续。}{A continuous function on a compact set is uniformly continuous: the same $\delta$ works at every point. This helps control approximation errors uniformly rather than requiring smaller and smaller steps at different locations. Continuity on a noncompact domain need not be uniform; $x^2$ is not uniformly continuous on the real line but is uniformly continuous on every closed bounded interval.}
\BOTH{一致连续的量词顺序是：对每个 $\varepsilon>0$，存在 $\delta>0$，使任意 $x,y$ 都满足 $\|x-y\|<\delta\Rightarrow|f(x)-f(y)|<\varepsilon$。普通连续允许 $\delta$ 依赖基点；一致连续要求同一个 $\delta$ 同时适用于定义域内所有基点。}{Uniform continuity means that for each $\varepsilon>0$, there is one $\delta>0$ such that all pairs $x,y$ in the domain satisfy $|f(x)-f(y)|<\varepsilon$ whenever $\|x-y\|<\delta$. Ordinary continuity allows $\delta$ to depend on the base point; uniform continuity uses the same $\delta$ everywhere.}
\EX{函数 $f(x,y)=x+y$ 在 $\mathbb R^2$ 一致连续，因为 $|f(x_1,y_1)-f(x_2,y_2)|\le\sqrt2\|(x_1,y_1)-(x_2,y_2)\|$。这条全局 Lipschitz 估计直接给出 $\delta=\varepsilon/\sqrt2$。}{The function $f(x,y)=x+y$ is uniformly continuous on $\mathbb R^2$, since its output difference is at most $\sqrt2$ times the Euclidean distance. This global Lipschitz bound gives $\delta=\varepsilon/\sqrt2$.}
\ASK{为什么检查多元函数沿几条直线不能证明极限存在？}{Why can checking a few lines not prove a multivariable limit?}{直线只覆盖所有趋近方向中的一部分；曲线路径的斜率还可随距离变化。证明需要控制整个去心邻域，可用范数不等式或序列判据。}{A few lines cover only some approaches, and curved paths can change direction with distance. A proof must control the entire punctured neighborhood, for example with a norm estimate or the sequential criterion.}

\section{级数判别与多元极限 | Series tests and limits}
\subsection{级数判别：先识别项，再检查临界条件}
\BOTH{判断 $\sum a_n$ 时，先检查 $a_n\to0$；否则级数立即发散。若 $a_n\ge0$，寻找比较对象或积分；含阶乘、指数时试比值；含 $n$ 次幂时试根值；交错符号时检查 Leibniz 条件。判别法有具体假设，不能只凭公式外形选择。}{For a series $\sum a_n$, first check whether $a_n\to0$; otherwise it diverges immediately. For nonnegative terms, look for a comparison or integral test. Factorials and exponentials suggest the ratio test; $n$th powers suggest the root test; alternating signs suggest Leibniz's test. Each test has hypotheses, so its formula alone does not determine when it applies.}
\BOTH{例一，$\sum 3^n/n!$：相邻项之比为 $3/(n+1)\to0<1$，故收敛。例二，$\sum n/(n^2+1)$：与 $1/n$ 极限比较，因为 $\frac{n/(n^2+1)}{1/n}=n^2/(n^2+1)\to1$，故发散。例三，$\sum(-1)^{n-1}/\sqrt n$：正项 $1/\sqrt n$ 递减趋零，所以交错级数收敛；绝对值级数是 $p=1/2$ 级数，发散，故原级数条件收敛。}{For $\sum 3^n/n!$, the ratio of consecutive terms tends to zero, so it converges. The terms $n/(n^2+1)$ are asymptotic to $1/n$, so their series diverges by limit comparison. The alternating series with terms $1/\sqrt n$ converges by Leibniz's test, but its absolute-value series diverges because it is a $p$-series with $p=1/2$; hence convergence is conditional.}
\BOTH{这些例子也说明判别结论的边界：比值极限小于 $1$ 能证明收敛，大于 $1$ 能证明发散，等于 $1$ 时没有结论；Leibniz 判别给出收敛，但若要知道是否绝对收敛，还须检查 $\sum|a_n|$。交错级数余项不超过首个遗漏项，所以第三个级数取前 $N$ 项的误差不超过 $1/\sqrt{N+1}$。}{These examples show the limits of the tests. A ratio limit below one proves convergence, above one proves divergence, and equal to one is inconclusive. Leibniz proves convergence but not absolute convergence; that requires checking $\sum|a_n|$. The alternating-series remainder is at most the first omitted term, so the error after $N$ terms in the third example is at most $1/\sqrt{N+1}$.}
\READ{“The ratio of consecutive terms tends to zero, which is less than one; therefore the series converges.”}

\subsection{用范数夹逼证明联合极限}
\BOTH{考虑 $f(x,y)=\frac{x^2y}{x^2+y^2}$，在原点以外定义，求 $(x,y)\to(0,0)$ 时的极限。沿直线逐条代入只能提供线索；证明需要对所有方向都成立的估计。因为 $0\le x^2/(x^2+y^2)\le1$，所以
\[
\left|\frac{x^2y}{x^2+y^2}\right|
=|y|\frac{x^2}{x^2+y^2}\le |y|
\le\sqrt{x^2+y^2}.
\]
右侧趋于零，由夹逼定理，联合极限为 $0$。}{Consider $f(x,y)=x^2y/(x^2+y^2)$ away from the origin. Checking individual lines only gives clues; a proof needs an estimate valid in every direction. Since $0\le x^2/(x^2+y^2)\le1$, the absolute value of $f$ is at most $|y|$, and $|y|\le\sqrt{x^2+y^2}$. The upper bound tends to zero, so the squeeze theorem proves that the joint limit is zero.}
\READ{“The absolute value of f of x comma y is at most the square root of x squared plus y squared.”}
\BOTH{对照函数 $g(x,y)=\frac{xy}{x^2+y^2}$：沿 $y=x$ 恒为 $1/2$，沿 $y=-x$ 恒为 $-1/2$，所以联合极限不存在。它的两个累次极限仍都为 $0$：先固定一个非零变量再令另一个趋零，函数趋零；在坐标轴上函数也为零。因此两个累次极限相等不能替代联合极限证明。}{By contrast, $g(x,y)=xy/(x^2+y^2)$ equals $1/2$ along $y=x$ and $-1/2$ along $y=-x$, so its joint limit does not exist. Its two iterated limits are nevertheless both zero: fix one nonzero coordinate and let the other tend to zero; on either axis the function is zero. Equal iterated limits therefore do not prove the joint limit.}
\ASK{多元极限的证明与反证分别怎么做？}{How do you prove or disprove a multivariable limit?}{反证时找两条趋近路径，使函数值趋向不同结果；证明时对所有趋近方式给出统一估计，例如 $|f(x)-L|\le C\|x-a\|^\alpha$，其中 $\alpha>0$。}{To disprove a limit, find two paths with different limiting values. To prove one, give a direction-independent estimate such as $|f(x)-L|\le C\|x-a\|^\alpha$ with $\alpha>0$.}

\subsection{紧致性如何保证实际问题有最优解}
\BOTH{设可行区域是闭单位圆盘 $K=\{(x,y):x^2+y^2\le1\}$，目标函数 $F(x,y)=3x+4y$ 表示线性收益。$K$ 闭且有界，所以在 $\mathbb R^2$ 中紧致；$F$ 连续，故最大值一定取得。Cauchy--Schwarz 不等式给出 $3x+4y\le5\sqrt{x^2+y^2}\le5$，在 $(x,y)=(3/5,4/5)$ 取等号，因此最大收益为 $5$。}{Let the feasible region be the closed unit disk $K=\{(x,y):x^2+y^2\le1\}$ and the objective be $F(x,y)=3x+4y$. The set is closed and bounded, hence compact in $\mathbb R^2$. Since $F$ is continuous, it attains a maximum. Cauchy--Schwarz gives $3x+4y\le5\sqrt{x^2+y^2}\le5$, with equality at $(3/5,4/5)$, so the maximum is five.}
\BOTH{紧致性保证最优值不会只在极限处逼近却永远取不到，不等式负责找出具体最优点。若可行集不闭或不有界，连续目标函数也可能没有最大值：$F(x)=x$ 在开区间 $(0,1)$ 上的上确界是 $1$，但没有最大点。}{Compactness guarantees that the optimum is attained rather than merely approached; the inequality identifies the optimizer. If a feasible set is not closed or bounded, even a continuous objective may fail to attain its supremum: $F(x)=x$ on $(0,1)$ has supremum one but no maximum.}
\section{完备性如何把 Cauchy 列变成实数 | Completeness and a Cauchy sequence}
\noindent 将 $\sqrt2$ 的十进位截断写成 $q_n=\lfloor10^n\sqrt2\rfloor/10^n$。每个 $q_n$ 都是有理数，且 $0\le\sqrt2-q_n<10^{-n}$。若 $m>n$，则
\[
|q_m-q_n|\le|q_m-\sqrt2|+|q_n-\sqrt2|<10^{-m}+10^{-n}<2\cdot10^{-n},
\]
因此 $(q_n)$ 是 Cauchy 列。它在 $\mathbb R$ 中收敛到 $\sqrt2$，却不可能在 $\mathbb Q$ 中有极限：若有有理极限 $q$，由极限唯一性必有 $q=\sqrt2$，矛盾。这个例子具体说明有理数不完备，而实数完备性填补了有理数中的“缺口”。
\BOTH{令 $q_n=\lfloor10^n\sqrt2\rfloor/10^n$，即取 $\sqrt2$ 的前 $n$ 位小数截断。它是有理数，且 $0\le\sqrt2-q_n<10^{-n}$，所以 $(q_n)$ 是 Cauchy 列。在实数中它收敛到 $\sqrt2$；但在有理数中没有极限，否则极限唯一性会迫使有理数等于 $\sqrt2$。这说明 $\mathbb Q$ 不完备，而 $\mathbb R$ 的完备性填补这些缺口。}{Let $q_n=\lfloor10^n\sqrt2\rfloor/10^n$, the decimal truncation after $n$ places. These rational numbers satisfy $0\le\sqrt2-q_n<10^{-n}$, so they form a Cauchy sequence. They converge to $\sqrt2$ in $\mathbb R$ but have no limit in $\mathbb Q$; otherwise uniqueness of limits would make a rational number equal to $\sqrt2$. Thus $\mathbb Q$ is incomplete, while completeness of $\mathbb R$ fills this gap.}
\EX{例如 $q_1=1.4$、$q_2=1.41$、$q_3=1.414$。这些有理数越来越接近，但无论继续取多少位有限小数，其极限仍是无理数 $\sqrt2$。}{The first terms are $1.4$, $1.41$, and $1.414$. These rationals get closer, but their limit remains the irrational number $\sqrt2$.}
\ASK{Cauchy 列在 $\mathbb Q$ 中为什么不一定收敛？}{Why need not every Cauchy sequence in $\mathbb Q$ converge in $\mathbb Q$?}{Cauchy 条件只要求各项彼此越来越接近，没有保证候选极限仍属于所选空间。上述有理数列在 $\mathbb R$ 中趋于 $\sqrt2$，而 $\sqrt2\notin\mathbb Q$。}{The Cauchy condition says that sequence terms become arbitrarily close, but does not ensure that the limit belongs to the chosen space. The rational sequence above tends to $\sqrt2$ in $\mathbb R$, and $\sqrt2$ is not rational.}

\end{document}
